\section{REST para registrar: refuerzo académico en SGRA}
\begin{frame}{Registro de solicitudes en SGRA}
\tema{Khriz Coronel · REST para crear un registro}
\begin{columns}[T,onlytextwidth]
\begin{column}{.48\textwidth}
{\Large El estudiante solicita\par refuerzo en una asignatura.}\par\medskip
La solicitud queda pendiente. El docente decide después la programación.
\end{column}
\begin{column}{.46\textwidth}
\textbf{Transparencia de acceso}\par\smallskip
El cliente usa la API sin conocer las tablas.\par\medskip
\textbf{Límite visible}\par\smallskip
Sin respuesta, el cliente no sabe si se guardó.
\end{column}
\end{columns}
\vspace{6mm}
\begin{center}
\begin{tikzpicture}[>=Stealth,box/.style={draw=borde,fill=white,align=center,minimum height=.8cm,font=\small}]
\node[box,text width=2.7cm] (c) at (0,0) {Cliente externo};
\node[box,text width=3.5cm,fill=suave] (a) at (4.5,0) {API REST\\Valida los campos};
\node[box,text width=2.8cm] (d) at (9.2,0) {Solicitud\\Pendiente};
\draw[->,acento,thick] (c)--node[above,font=\scriptsize]{POST + JSON}(a);
\draw[->,acento,thick] (a)--node[above,font=\scriptsize]{Registrar}(d);
\end{tikzpicture}
\end{center}
\kgnota{Caso basado en RF01 del SGRA~\cite{sgra-rf01-khriz}. Demo autónoma con datos ficticios.}
\note{60 segundos. Presentar el Sistema de Gestión de Refuerzo Académico. Como ejemplo conocido, el estudiante solicita apoyo al docente; el contexto documentado recoge WhatsApp y correo como medios del proceso previo. Distinguir el proceso institucional, el backend anterior y esta demo. El horario indicado expresa disponibilidad; no es una sesión confirmada. La transparencia evita conocer las tablas, pero el fallo de comunicación debe mostrarse.}
\end{frame}

\begin{frame}[fragile]{Creación y rechazo mediante POST}
\tema{POST /api/student/requests}
\begin{columns}[T,onlytextwidth]
\begin{column}{.57\textwidth}
\small
\begin{verbatim}
{
  "subjectId": 1,
  "sessionTypeId": 1,
  "reason": "Necesito reforzar redes.",
  "preferredSchedule": "Martes 15:00"
}
\end{verbatim}
\end{column}
\begin{column}{.38\textwidth}
{\color{acento}\Large\bfseries 201 Created}\par\smallskip
\verb|{"requestId":1}|\par\medskip
{\color{ambar}\Large\bfseries 400 Bad Request}\par\smallskip
Sin motivo: corregir el campo. No se guarda un registro.
\end{column}
\end{columns}
\vspace{3mm}
\kgnota{La llamada espera una respuesta. Mensajería sería una alternativa con otro acuse y seguimiento~\cite{kazanavicius-comunicacion-2023-khriz}.}
\kgnota{Código y pruebas reproducibles en \texttt{demos/khriz/}. Registros en memoria durante la ejecución.}
\note{45 segundos. Mostrar los archivos de respuesta de la ejecución conservada, o ejecutar los dos curl de la guía. El primer identificador es 1 al iniciar una instancia nueva. Si falta reason, HTTP 400 identifica el campo y no inserta. La demo usa JSON; el backend anterior usa multipart y responde 200. El prototipo se adapta a la ficha de EXP-1. HTTP permite un intercambio directo; una cola necesitaría seguimiento del resultado. El artículo compara alternativas, no demuestra el rendimiento del SGRA.}
\end{frame}

\begin{frame}{Reintentos tras perder una respuesta}
\tema{Propuesta de idempotencia}
\begin{enumerate}
\item La API guarda la solicitud.
\item La respuesta se pierde y el cliente reintenta.
\item Sin control, aparece otra solicitud.
\end{enumerate}
\vspace{3mm}
\textbf{Una clave por operación}\par\smallskip
El reintento conserva la clave. El servicio recuperaría el identificador original.
\vspace{4mm}
\kgnota{La clave, el contenido y el resultado se guardarían junto con el registro, de forma atómica y duradera.}
\kgnota{En esta demo, repetir un POST válido crea otro registro. La idempotencia queda propuesta.}
\note{30 segundos. El timeout no prueba que el servidor no guardó. La clave se crea para una operación y se conserva al reintentarla; una operación nueva recibe otra clave. En una integración real se relaciona con la identidad y el contenido, y se conserva atómicamente con el resultado. Reutilizar una clave para otro contenido debe rechazarse. No confundirlo con impedir solicitudes distintas que coinciden en una asignatura.}
\end{frame}

\begin{frame}{Coordinación de un último cupo}
\tema{CAP y PACELC · arquitectura futura}
\begin{center}
\begin{tikzpicture}[>=Stealth,box/.style={draw=borde,fill=white,align=center,text width=3.6cm,minimum height=.9cm,font=\small}]
\node[box] (a) at (0,0) {Nodo A\\Confirma el cupo};
\node[box] (b) at (7,0) {Nodo B\\Confirma el mismo cupo};
\draw[<->,acento,thick] (a)--node[above,font=\scriptsize]{Coordinación}(b);
\draw[ambar,very thick] (3.3,-.22)--(3.7,.22);
\draw[ambar,very thick] (3.3,.22)--(3.7,-.22);
\end{tikzpicture}
\end{center}
\begin{columns}[T,onlytextwidth]
\begin{column}{.47\textwidth}
\textbf{Con partición · CAP}\par\smallskip
Si no se puede garantizar el cupo único, suspender su confirmación~\cite{gilbert-lynch-cap-khriz}.
\end{column}
\begin{column}{.47\textwidth}
\textbf{Sin partición · PACELC}\par\smallskip
Esperar coordinación preserva la decisión consistente y añade latencia~\cite{abadi-pacelc-khriz}.
\end{column}
\end{columns}
\vspace{4mm}
\kgnota{Registrar una solicitud pendiente no confirma un cupo. La demo no implementa réplicas ni mide estos compromisos.}
\note{45 segundos. Señalar dos nodos que gestionan el mismo cupo y pierden la comunicación. La regla de la propuesta evita dos confirmaciones incompatibles; se rechaza o pospone confirmar cuando falta coordinación suficiente. La pérdida de un enlace no obliga a detener todo el sistema. PACELC considera, sin partición, la latencia de coordinar para mantener la consistencia. Estos compromisos pertenecen a la arquitectura propuesta, no al uso de REST ni a la idempotencia.}
\end{frame}
